\subsection{Cinemática inversa}
Se sigue la reducción geométrica de las semanas 5 y 7 \parencite{course5,course7}. Para un objetivo $(x,y,z)$, la distancia horizontal es $r=\sqrt{x^2+y^2}$ y la altura respecto del hombro es $h=z-l_1$. Fuera del eje, con $\phi=\operatorname{atan2}(y,x)$, dos pares de ángulo de base y proyección radial satisfacen $x=\rho\cos\theta_1$, $y=\rho\sin\theta_1$:
\begin{equation}
(\theta_1,\rho)=(\phi,+r)\quad\text{o}\quad(\phi+\pi,-r).
\end{equation}
El triángulo existe si $|l_2-l_3|\leq D\leq l_2+l_3$, con $D^2=r^2+h^2$. Para cada base se calculan el codo y el hombro mediante
\begin{align}
C_3&=\frac{D^2-l_2^2-l_3^2}{2l_2l_3},\qquad \theta_3=\pm\arccos C_3,\\
\theta_2&=\operatorname{atan2}(h,\rho)-
\operatorname{atan2}(l_3\sin\theta_3,l_2+l_3 C_3).
\end{align}
Para el punto $(0.333420,0.192500,0.291147)$ m, con coordenadas redondeadas, se obtienen cuatro ternas en grados:
\begin{center}
\begin{tabular}{rrrr}\toprule
$\rho$ (m) & Base & Hombro & Codo\\\midrule
$+0.385$ & 30 & 6.24 & 60\\
$+0.385$ & 30 & 60 & $-60$\\
$-0.385$ & $-150$ & 120 & 60\\
$-0.385$ & $-150$ & 173.76 & $-60$\\\bottomrule
\end{tabular}
\end{center}
Cada terna se verifica a precisión completa con las ecuaciones trigonométricas de posición y la traslación de $M_{01}M_{12}M_{23}$. En extensión o plegado coinciden las ramas de codo. Sobre el eje, cualquier base alcanza la misma posición y se elige un representante. La viabilidad física de cada postura requiere conocer los límites articulares y las colisiones.
